\vspace{-16pt}
\section{Introduction}
Iterative refinement is a natural way to use an instruction-based image editor. A user might ask the model to make the sky look like sunset, then add a bird, then remove a fence in the background, refining an image through several small edits rather than one large prompt. This pattern introduces a failure mode that single-step evaluation cannot see: a model can follow every individual instruction correctly while still drifting further and further from what the user actually wanted, simply because nothing in its training or inference constrains it to leave the rest of the image alone. Existing evaluation protocols for these models, including CLIP-based instruction-following scores and human preference studies, measure whether the requested change happened. None of them measure what else changed as a side effect. This paper calls that side effect \textit{semantic drift}.

Most existing instruction-based editors, InstructPix2Pix chief among them \cite{ref_ip2p}, are trained on synthetically generated (image, instruction, edited-image) triplets. That training recipe supplies no explicit signal for "leave everything else untouched" beyond whatever the data-generation process happened to preserve. Drift is therefore an expected consequence of how these models are built, not merely an empirical curiosity worth measuring after the fact.

Three questions follow directly from this observation. Can unintended drift be measured automatically, without manual inspection of every edited image (RQ1)? Does drift accumulate as a chain of edits gets longer, so that later instructions cause more collateral damage than earlier ones (RQ2)? And can cheap, pretrained-model-only interventions reduce it (RQ3)? This paper answers all three.

The contribution is threefold. First, the Drift Score, a reusable metric combining Segment Anything (SAM) region segmentation with CLIP embedding similarity to quantify unintended regional change, at both the single-edit and full-chain level. Second, an empirical study of how drift behaves across 120 real four-step edit chains, split between localized object-level edits and broad scene-level edits. Third, a statistically grounded comparison of two mitigation strategies, masked conditioning and region-locking, against an unmitigated baseline.

The work was scoped for a single contributor on a 30-day timeline using only pretrained models and free-tier compute: a laptop for development and analysis, and a free Google Colab GPU for the heavier inference stages. No model was trained or fine-tuned. The contribution here is the measurement methodology and the empirical findings it produces, not a new editing model.
